\documentclass[aspectratio=169]{beamer}
\usepackage{tikz}
\usepackage{graphicx}
\usepackage{listings}
\usepackage{xcolor}
\usepackage{hyperref}
\usepackage[spanish,es-noquoting]{babel}
\usepackage[round]{natbib}

\lstset{
  basicstyle=\ttfamily\small,
  keywordstyle=\color{blue},
  stringstyle=\color{red},
  commentstyle=\color{gray},
  showstringspaces=false
}

% Clean Theme Setup & Optimization
\usetheme{Warsaw}
\usecolortheme{default}
\usefonttheme{serif}

% Remove top headline bar and navigation clutter to maximize slide text space
\setbeamertemplate{headline}{}
\setbeamertemplate{navigation symbols}{}

% Slim, clean footline bar
\setbeamertemplate{footline}{
  \leavevmode%
  \hbox{%
    \begin{beamercolorbox}[wd=.333333\paperwidth,ht=2.25ex,dp=1ex,left]{author in head/foot}%
      \hspace*{2ex}\insertshortauthor
    \end{beamercolorbox}%
    \begin{beamercolorbox}[wd=.333333\paperwidth,ht=2.25ex,dp=1ex,center]{title in head/foot}%
      \insertshorttitle
    \end{beamercolorbox}%
    \begin{beamercolorbox}[wd=.333333\paperwidth,ht=2.25ex,dp=1ex,right]{date in head/foot}%
      \insertframenumber{} / \inserttotalframenumber\hspace*{2ex}
    \end{beamercolorbox}%
  }%
  \vskip0pt%
}

% Font size and list spacing optimizations to prevent text overflow
\setbeamertemplate{itemize/enumerate body begin}{\small}
\setbeamertemplate{itemize/enumerate subbody begin}{\footnotesize}
\setbeamertemplate{itemize/enumerate subsubbody begin}{\scriptsize}

\makeatletter
\g@addto@macro\itemize{%
  \setlength{\itemsep}{2pt}%
  \setlength{\parsep}{1pt}%
  \setlength{\topsep}{2pt}%
  \setlength{\partopsep}{0pt}%
}
\makeatother

\graphicspath{{../}{./}}
\setbeamertemplate{caption}[numbered]

% Title info
\title[LegalAI Presentation]{LegalAI Presentation}
\subtitle{}
\author{Ricardo Reyes-Jimenez}
\institute{Universidad de La Sabana}
\date{\today}



\begin{document}

% Title slide
\begin{frame}
  \titlepage
\end{frame}

% Reset background for normal slides
\usebackgroundtemplate{}


\begin{frame}[allowframebreaks,fragile]{Automatización de Tareas Usando Inteligencia Artificial}

\textbf{Una Perspectiva del Analista de Datos Senior}
\textbf{[Tu Nombre/Experto]}
\textbf{[Tu Institución/Organización]}

\end{frame}

\begin{frame}[allowframebreaks,fragile]{Objetivos de la Sesión y Contexto General}


\textbf{Objetivos}
\begin{itemize}
  \item Comprender el impacto transformador de la IA en la automatización de tareas de análisis de datos.
  \item Identificar oportunidades y desafíos clave para el Analista de Datos Senior.
  \item Explorar casos de estudio y ejemplos prácticos de automatización inteligente.
  \item Discutir las implicaciones éticas y la protección de datos en sistemas automatizados.
\end{itemize}
\textbf{Contexto General}
La Inteligencia Artificial está redefiniendo el panorama del análisis de datos, pasando de la ejecución de tareas repetitivas a la toma de decisiones inteligentes y la optimización de flujos de trabajo completos. Para el Analista de Datos Senior, es crucial dominar esta evolución.

\end{frame}

\begin{frame}[allowframebreaks,fragile]{Marco Conceptual y Conceptos Clave}

\textbf{¿Qué es la Automatización con IA en Análisis de Datos?}
La automatización con IA va más allá de la simple programación de scripts. Implica sistemas capaces de:
\begin{itemize}
  \item \textbf{Aprender y adaptarse:} Mejorar su rendimiento con base en nuevos datos y experiencias.
  \item \textbf{Reconocer patrones complejos:} Identificar tendencias y anomalías en grandes volúmenes de datos.
  \item \textbf{Tomar decisiones:} Ejecutar acciones o recomendar cursos de acción basados en el análisis.
\end{itemize}
\end{frame}

\begin{frame}[allowframebreaks,fragile]{Marco Conceptual y Conceptos Clave}

\textbf{Conceptos Clave}
\begin{itemize}
  \item \textbf{Agentes Autónomos de IA:} Sistemas capaces de gestionar y ejecutar tareas de análisis de datos de principio a fin (end-to-end), desde la ingesta de datos hasta la generación de informes (Javad et al., 2026).
  \item \textbf{Procesamiento de Datos en Tiempo Real:} La capacidad de procesar y analizar grandes volúmenes de datos a medida que se generan, permitiendo insights y decisiones instantáneas (Seethala, 2025).
  \item \textbf{Análisis Multi-Ómico y Datos Complejos:} Automatización para extraer patrones y conocimientos de conjuntos de datos altamente complejos y voluminosos (Oskolkov, 2025).
  \item \textbf{Decisiones Automatizadas:} Sistemas de IA que toman decisiones directamente o asisten en ellas, lo que implica consideraciones críticas sobre ética y protección de datos (Unknown Authors, 2022).
\end{itemize}
\end{frame}

\begin{frame}[allowframebreaks,fragile]{Casos de Estudio y Ejemplos Prácticos Aplicados}

\textbf{1. Automatización Integral de Flujos de Análisis}
\begin{itemize}
  \item \textbf{Descripción:} Implementación de agentes de IA para gestionar todo el ciclo de vida del análisis de datos:
  \begin{itemize}
    \item \textbf{Ingesta y Limpieza de Datos:} Identificación y corrección automática de errores, estandarización de formatos.
    \item \textbf{Transformación y Modelado:} Selección automática de características, ejecución de modelos predictivos y descriptivos.
    \item \textbf{Generación de Informes:} Creación automática de dashboards y resúmenes con insights clave.
  \end{itemize}
  \item \textbf{Beneficio:} Reducción significativa del tiempo y esfuerzo manual, permitiendo a los analistas enfocarse en la interpretación y estrategia.
  \item \textbf{Referencia:} Javad et al., 2026.
\end{itemize}
\end{frame}

\begin{frame}[allowframebreaks,fragile]{Casos de Estudio y Ejemplos Prácticos Aplicados}

\textbf{2. Optimización de Almacenamiento y Análisis en Salud}
\begin{itemize}
  \item \textbf{Descripción:} Uso de IA para automatizar el almacenamiento y análisis de datos de pacientes en tiempo real dentro de entornos de salud:
  \begin{itemize}
    \item \textbf{Monitorización Continua:} Procesamiento instantáneo de datos de sensores, dispositivos médicos y registros electrónicos.
    \item \textbf{Detección de Anomalías:} Identificación temprana de riesgos para la salud o cambios críticos en el estado del paciente.
    \item \textbf{Soporte a la Decisión Clínica:} Provisión de insights inmediatos a profesionales médicos para tratamientos proactivos.
  \end{itemize}
  \item \textbf{Beneficio:} Mejora de la calidad de la atención al paciente y eficiencia operativa.
  \item \textbf{Referencia:} Seethala, 2025.
\end{itemize}
\end{frame}

\begin{frame}[allowframebreaks,fragile]{Ejercicios Prácticos y Actividades Interactivas}

\textbf{Ejercicio 1: Identificando Oportunidades de Automatización en Tu Rol}
\begin{itemize}
  \item \textbf{Instrucción:} Piensa en una tarea recurrente en tu trabajo actual como analista de datos.
  \item \textbf{Pregunta Guía:}
  \begin{itemize}
    \item ¿Qué pasos manuales consumes más tiempo en esa tarea?
    \item ¿Cómo podría la IA automatizar o asistir en esos pasos (ej. limpieza de datos, generación de informes, detección de patrones)?
    \item ¿Qué beneficios específicos esperas obtener de esa automatización?
  \end{itemize}
  \item \textbf{Actividad:} Comparte tu tarea y una idea de automatización con un compañero.
\end{itemize}
\end{frame}

\begin{frame}[allowframebreaks,fragile]{Ejercicios Prácticos y Actividades Interactivas}

\textbf{Ejercicio 2: Dilemas Éticos en Decisiones Automatizadas}
\begin{itemize}
  \item \textbf{Contexto:} Imagina un sistema de IA que automatiza la decisión de otorgar un crédito o una póliza de seguro, basándose en el historial de datos.
  \item \textbf{Preguntas Guía:}
  \begin{itemize}
    \item ¿Qué riesgos éticos o de privacidad podrían surgir si la IA toma la decisión final sin supervisión humana?
    \item ¿Cómo asegurarías la equidad y transparencia del algoritmo para evitar sesgos involuntarios?
    \item ¿Qué mecanismos de supervisión y auditabilidad deberían implementarse?
  \end{itemize}
  \item \textbf{Actividad:} Discute en grupos pequeños cómo mitigar estos riesgos.
\end{itemize}
\end{frame}

\begin{frame}[allowframebreaks,fragile]{Sabías que... Datos Curiosos}

\begin{itemize}
  \item \textbf{El Rol del Analista Cambia de 'Hacer' a 'Supervisar':} La emergencia de agentes de IA autónomos que gestionan el análisis de datos de principio a fin significa que el Analista de Datos Senior transitará de la ejecución manual a la supervisión estratégica, la validación de modelos y la gobernanza ética de la IA. ¡Esto requiere una profunda comprensión de las limitaciones y sesgos de los modelos!
  \item \textbf{Valor en Tiempo Real Transforma Industrias:} La capacidad de la automatización con IA para procesar datos de pacientes en tiempo real, especialmente en entornos críticos como la atención médica, permite una generación de insights continua e instantánea. Este cambio de procesamiento por lotes a respuestas dinámicas redefine fundamentalmente cómo los datos pueden impactar las decisiones operativas y estratégicas.
  \item \textbf{La Ética es el Corazón, No un Apéndice:} A medida que la IA automatiza más `\texttt{decisiones'', la protección de datos y las consideraciones éticas no son un }\texttt{extra'', sino un principio de diseño central. Los Analistas de Datos Senior deben dominar las implicaciones legales, éticas y sociales, garantizando la equidad, transparencia y responsabilidad en los sistemas impulsados por IA. La }`ética de datos'' es ahora una competencia esencial.
\end{itemize}
\end{frame}

\begin{frame}[allowframebreaks,fragile]{Conclusiones y Puntos Clave (Takeaways)}

\begin{itemize}
  \item \textbf{Transformación del Rol del Analista:} La automatización con IA impulsa al Analista de Datos Senior hacia un rol más estratégico, centrado en la supervisión, validación, interpretación y gobernanza ética de los sistemas de IA.
  \item \textbf{Eficiencia Operacional sin Precedentes:} La IA permite la automatización de flujos de trabajo de análisis de principio a fin y el procesamiento de datos en tiempo real, lo que se traduce en una eficiencia y una capacidad de respuesta inigualables.
  \item \textbf{La Ética y la Transparencia son Fundamentales:} A medida que las decisiones se automatizan, la protección de datos, la equidad algorítmica y la transparencia se convierten en pilares esenciales para el diseño y la implementación de soluciones de IA.
  \item \textbf{Oportunidades de Crecimiento:} Para el Analista de Datos Senior, dominar estas tendencias no es solo una ventaja, sino una necesidad para liderar la innovación y la transformación digital en sus organizaciones.
\end{itemize}
\end{frame}

\begin{frame}[allowframebreaks,fragile]{Referencias Bibliográficas}


\begin{itemize}
  \item Javad, H., Kc, P., \& Zhgair, H. (2026). \textit{Toward Integrated AI-Powered Autonomous Agents for End-to-End Data Analysis and Task Automation}. Proceedings of the International Conference on Industrial Engineering and Operations Management. DOI: 10.46254/na11.20260260
  \item Seethala, S. (2025). \textit{Next-Level Data Warehousing in Healthcare: AI-Powered Automation for Real Time Patient Data}. Unknown Journal. DOI: 10.2139/ssrn.5113306
  \item Oskolkov, N. (2025). \textit{AI-Powered multiOmics Data Analysis}. Unknown Journal. DOI: 10.61700/vkpkhmnwa12rv2031
  \item Unknown Authors. (2022). \textit{La automatización de las decisiones}. Decisiones automatizadas y protección de datos. Especial atención a los sistemas de inteligencia artificial. DOI: 10.2307/j.ctv2zp4t7t.7
\end{itemize}
\end{frame}




\end{document}